\documentclass[../main.tex]{subfiles}
\begin{document}

\section{Интерфейс системы}

\subsection{Общая структура главного окна}

Интерфейс выполнен в виде одностраничного приложения. Все элементы
расположены в одном окне и не требуют перехода между отдельными
страницами. Главное окно разделено на левую панель источников
и рабочую область. Такое расположение позволяет одновременно видеть
список подключённых баз, структуру выбранной таблицы и результат
её просмотра или сравнения.

Структура главного окна показана на рисунке~\ref{fig:ui-layout}.
В левой части находится форма подключения нового источника, список
загруженных инстансов и дерево «инстанс~--- схема~--- таблица».
В правой части расположены переключатель режимов «Данные» и «Джойн»,
рабочая область и выдвижная панель сравнения схем.

\begin{figure}[H]
  \centering
  \resizebox{\linewidth}{!}{%
    \input{\subfix{diagrams/uiLayout}}%
  }
  \caption{Структура главного окна}
  \label{fig:ui-layout}
\end{figure}

Для каждого инстанса доступна кнопка обновления схемы. Форма
подключения содержит название источника, JDBC-URL,
имя пользователя, пароль, необязательный alias и тип СУБД. После
создания источник появляется в общем списке, а его схема
загружается в дерево. На рисунке~\ref{fig:ui-main} показано главное
окно с формой подключения и списком источников.

\screenshot{proofs/ui-main.png}
  {Главное окно: подключение источника и список инстансов}
  {fig:ui-main}
  {Создание инстанса, список из трёх источников, дерево схем.}

\subsection{Работа с деревом схемы}

Дерево схемы отображает иерархию объектов источника. При раскрытии
таблицы показываются её колонки, ограничения и индексы. Для каждой
колонки указаны тип, признак \field{NOT NULL}, значение по умолчанию
и признак генерации. Ограничения отображаются с составом ключа
и правилами \field{ON DELETE}/\field{ON UPDATE}; индексы~--- с типом
и перечнем колонок. Это позволяет изучить структуру таблицы, не
открывая отдельный клиент СУБД.

При выборе таблицы в дереве в рабочей области открывается панель
данных (рисунок~\ref{fig:ui-schema}). Выбранная таблица выделяется
в дереве, а рядом с ней показывается число колонок, ограничений
и индексов.

\screenshot{proofs/ui-schema.png}
  {Дерево схемы: колонки, ограничения и индексы}
  {fig:ui-schema}
  {Раскрытая таблица с колонками, первичным ключом, внешними ключами
   и индексами.}

\subsection{Просмотр и сравнение данных}

Панель данных выполняет постраничную выборку строк. Строки
упорядочиваются по первичному ключу, что обеспечивает
детерминированную пагинацию; размер страницы по умолчанию равен
100, максимальный~--- 200. В заголовке панели указываются имя
таблицы, число колонок и состав первичного ключа. На
рисунке~\ref{fig:ui-rows} приведён просмотр строк таблицы.

\screenshot{proofs/ui-rows.png}
  {Постраничный просмотр строк таблицы}
  {fig:ui-rows}
  {Таблица с большим числом строк, пагинация, отображение первичного
   ключа.}

Для контроля изменений данных используется построчное сравнение:
клиент запоминает строки, отображавшиеся на экране, читает свежие
строки источника и сравнивает их по первичному ключу. Добавленные,
удалённые и изменённые строки подсвечиваются, а для изменённых
показываются значения «было» и «стало». Результат приведён
на рисунке~\ref{fig:ui-row-diff}.

\screenshot{proofs/ui-row-diff.png}
  {Построчное сравнение данных: эталон и источник}
  {fig:ui-row-diff}
  {Кнопка «Обновить и сравнить», подсветка добавленных, удалённых
   и изменённых строк.}

\subsection{Сравнение схем}

Сравнение схемы можно запустить из панели источника. Сервис
собирает текущую схему источника и сравнивает её с сохранённым
каталогом по именам и сигнатурам объектов, а не по внутренним
идентификаторам. Результат открывается в выдвижной панели
\field{SchemaDiffDrawer} (рисунок~\ref{fig:ui-schema-diff}).
Добавленные, удалённые и изменённые объекты группируются по схемам
и таблицам, поэтому можно быстро оценить масштаб расхождений.

\screenshot{proofs/ui-schema-diff.png}
  {Сравнение сохранённой и текущей схемы}
  {fig:ui-schema-diff}
  {Панель сравнения схем с добавленными, удалёнными и изменёнными
   объектами.}

\subsection{Конструктор федеративных запросов}

В режиме «Джойн» открывается визуальный конструктор
(рисунок~\ref{fig:ui-join}). Пользователь добавляет таблицы из разных
источников, отмечает колонки результата и задаёт условия соединения.
Конструктор не требует писать SQL вручную: он формирует структуру
запроса и отправляет её в \path{POST /api/join}. Таблицы можно
добавлять из дерева схемы кнопкой «+» или через поиск по полному
имени \field{alias.schema.table}.

\screenshot{proofs/ui-join.png}
  {Конструктор федеративного запроса}
  {fig:ui-join}
  {Таблицы из разных источников, выбранные колонки и условие
   соединения.}

Результат запроса отображается в таблице с указанием количества
строк и длительности выполнения (рисунок~\ref{fig:ui-join-result}).
Отдельная вкладка «История» показывает журнал выполненных запросов:
SQL, статус, число строк и время выполнения
(рисунок~\ref{fig:ui-join-history}).

\screenshot{proofs/ui-join-result.png}
  {Результат федеративного запроса}
  {fig:ui-join-result}
  {Таблица результата, количество строк и длительность выполнения.}

\screenshot{proofs/ui-join-history.png}
  {История федеративных запросов}
  {fig:ui-join-history}
  {Журнал запросов: SQL, статус, число строк и длительность.}

\end{document}
